\subsection{Class-Builder Notation}
\label{sec:technical-class-builders}

The primitive language of NBG contains no terms other than variables.
Nevertheless, the extended syntax introduced in this text permits
expressions of the form
\[
\{\xi\mid\varphi\}
\]
to be used as though they were terms denoting classes. We now specify
this extension and establish the formal basis for its use.

The intended role of such an expression is to represent the class
whose members are precisely the sets satisfying \(\varphi\). This
requires two distinct justifications. First, every class-builder
expression must represent a uniquely determined class. Second, every
complete formula containing class-builder expressions must admit
elimination in favor of a formula of the primitive language.

There is a further requirement arising from the predicative restriction
on Class Formation. A formula containing class-builder expressions may
itself occur as the defining condition of another class-builder
expression. We must therefore establish that the introduction of
class-builder notation does not permit the formation of classes whose
existence is not justified by the original Class Formation Schema.

These requirements will be addressed through a construction organized
by the complexity of class-builder expressions. We begin by specifying
the expressions involved and the procedures by which they are connected
to the existing language.

\subsubsection{Class-Term Syntax and Contextual Elimination}

\nbgpara{1}{General Class Terms}

We extend the syntax by admitting a new category of expressions, called
\emph{extended terms}.

\begin{definition}[General Class Terms]
\label{def:general-class-terms}
Every variable is an extended term.

If \(\varphi\) is a predicative extended formula and \(\xi\) is a
variable, then
\[
\{\xi\mid\varphi\}
\]
is a \emph{general class term}, provided that every class term occurring
in \(\varphi\) has already been admitted at an earlier stage of the
construction.

Every general class term is an extended term.

If \(S\) and \(T\) are extended terms, then
\[
S\in T
\qquad\text{and}\qquad
S=T
\]
are atomic extended formulas.
\end{definition}

The ordinary formation rules for formulas extend to these atomic
extended formulas. Thus, extended formulas containing class terms may
be combined using the logical connectives and quantifiers, as well as
previously admitted defined notation.

The variable \(\xi\) in
\[
\{\xi\mid\varphi\}
\]
binds every occurrence of \(\xi\) free in \(\varphi\). The free
variables of the resulting class term are therefore precisely the
variables occurring free in \(\varphi\), other than \(\xi\).

The conventions concerning free and bound occurrences, alphabetic
variation, and substitution of variables for variables extend to
class terms in the corresponding recursive manner. In particular,
the bound variable of a class term may be renamed alphabetically
to avoid capture. The compatibility of these operations with
primitive expansion will be established below.

The formation clauses do not yet license the substitution of arbitrary
extended terms for variables. This more general operation will be
justified separately, once the eliminability of class terms has been
established.

\nbgpara{2}{Bounded Class Terms}

We also admit a bounded form of class-builder notation.

\begin{definition}[Bounded Class Terms]
\label{def:bounded-class-terms}
If \(\xi\) and \(\eta\) are distinct variables and \(\varphi\) is a
predicative extended formula whose class terms, if any, have already
been admitted at earlier stages, then
\[
\{\xi\in\eta\mid\varphi\}
\]
is a \emph{bounded class term}.

Bounded class terms are defined by the expansion
\[
\{\xi\in\eta\mid\varphi\}
\defequiv
\{\xi\mid(\xi\in\eta\land\varphi)\}.
\]
\end{definition}

The variable \(\xi\) binds its free occurrences in \(\varphi\),
but does not bind \(\eta\). Thus, the free variables of
\[
\{\xi\in\eta\mid\varphi\}
\]
are the variables free in \(\varphi\) other than \(\xi\), together
with \(\eta\).

Bounded class terms introduce no independent class-forming operation.
Every occurrence is understood through its prescribed expansion into
a general class term. In particular, the existence and uniqueness
of the class represented by a bounded term will follow from the
corresponding results for general class terms.

The restriction to variable bounds is provisional. Once substitution
of extended terms has been justified, bounded class-builder notation
will be extended to arbitrary admitted extended terms as bounds.

\nbgpara{3}{Builder Rank}

Class terms may occur within the defining conditions of further class
terms. To organize their introduction without circularity, we assign
a \emph{builder rank} to extended expressions.

\begin{definition}[Builder Rank]
\label{def:builder-rank}
Variables have builder rank \(0\). An extended formula containing no
class term also has builder rank \(0\).

If
\[
T=\{\xi\mid\varphi\},
\]
then
\[
\operatorname{br}(T)
=
\operatorname{br}(\varphi)+1.
\]

A bounded class term has the builder rank of the general class term
to which it is expanded. Thus,
\[
\operatorname{br}
\bigl(
\{\xi\in\eta\mid\varphi\}
\bigr)
=
\operatorname{br}(\varphi)+1.
\]

The builder rank of any other extended expression containing class
terms is the greatest builder rank of the class terms occurring
within it.
\end{definition}

The formation of class terms is understood inductively in builder
rank.

At rank \(0\), the extended terms are the variables, and the
predicative extended formulas are those admitted before the
introduction of class-builder notation.

Suppose that class terms of rank at most \(n\), together with the
predicative extended formulas containing them, have been admitted.
A new class term
\[
\{\xi\mid\varphi\}
\]
may then be introduced if \(\varphi\) is predicative according to
the reductions already justified at these earlier ranks. Its rank
is
\[
\operatorname{br}(\varphi)+1.
\]

Consequently, a class term of rank \(n+1\) has a defining formula
containing only class terms of rank at most \(n\).

Predicativity and builder rank impose different requirements.
Predicativity determines which defining conditions may be used
in applications of Class Formation. Builder rank orders the
justification of nested class terms, ensuring that the admissibility
of a new term depends only upon results already established for
terms of lower rank.

The fact that this construction can be carried out through every
finite rank will be established by the Class-Term Extension
Metatheorem.

\nbgpara{4}{Canonical Membership Conditions}

We associate with each class term a formula expressing the condition
under which a variable belongs to the class represented by that term.

\begin{definition}[Canonical Membership Condition]
\label{def:technical-canonical-membership-condition}
Let
\[
T=\{\xi\mid\varphi\}
\]
be a general class term, and let \(\zeta\) be a variable free for
\(\xi\) in \(\varphi\). The \emph{canonical membership condition}
of \(T\), with argument \(\zeta\), is
\[
\operatorname{Cond}_T(\zeta)
\defequiv
\Set(\zeta)\land\varphi[\zeta/\xi].
\]
\end{definition}

The sethood condition is essential. Under the intended interpretation
of NBG, every member of a class is a set. Thus, even if a proper class
satisfies the defining formula \(\varphi\), it cannot thereby become a
member of the class represented by \(T\).

For a bounded class term
\[
T=\{\xi\in\eta\mid\varphi\},
\]
the canonical membership condition is determined through its
expansion into a general class term:
\[
\operatorname{Cond}_T(\zeta)
\defequiv
\Set(\zeta)\land
\bigl(
\zeta\in\eta\land\varphi[\zeta/\xi]
\bigr).
\]
Since membership in \(\eta\) already implies sethood, this condition
is provably equivalent to
\[
\zeta\in\eta\land\varphi[\zeta/\xi].
\]
The resulting simplification will be used in the Membership Law for
bounded class terms.

The notation
\[
\operatorname{Cond}_T(\zeta)
\]
is metalinguistic. The expression ``\(\operatorname{Cond}\)''
is not a symbol of the primitive language or of the extended
syntax. It serves only to abbreviate the indicated extended formula.

Whenever the substitution occurring in a canonical membership
condition is not immediately legitimate, the bound variable of
the class term may first be renamed alphabetically. In particular,
a suitable canonical membership condition can always be formed
using a variable not occurring in the class term.

If \(T\) has builder rank \(n+1\), then
\(\operatorname{Cond}_T(\zeta)\) contains only class terms of rank
at most \(n\). This rank-lowering property will be fundamental
to the subsequent construction.

\nbgpara{5}{Realizations}

A class term is connected to the existing language by specifying
what it means for a variable to represent the class determined
by its canonical membership condition.

\begin{definition}[Realization of a Class Term]
\label{def:technical-realization-of-class-term}
Let \(T\) be a general class term, and let \(\eta\) be a variable.
Choose a variable \(\xi\) distinct from \(\eta\) and not occurring
in \(T\). Then
\[
\Real_T(\eta)
\defequiv
\forall\xi
\bigl(
\xi\in\eta
\leftrightarrow
\operatorname{Cond}_T(\xi)
\bigr).
\]
\end{definition}

The formula
\[
\Real_T(\eta)
\]
expresses that \(\eta\) has precisely the members specified by
the canonical membership condition of \(T\). When this formula
holds, we say that \(\eta\) \emph{realizes} \(T\), or that
\(\eta\) is a \emph{realization} of \(T\).

The notation \(\Real_T(\eta)\) is metalinguistic, just like
\(\operatorname{Cond}_T(\xi)\). It abbreviates an extended formula
and introduces no additional symbol into the primitive language.
Different legitimate choices of the bound variable \(\xi\)
yield alphabetic variants.

For bounded class terms, realization is defined through the
prescribed expansion into a general class term. Thus, if
\[
T=\{\xi\in\eta\mid\varphi\},
\]
then
\[
\Real_T(\zeta)
\]
abbreviates
\[
\Real_{\{\xi\mid(\xi\in\eta\land\varphi)\}}(\zeta).
\]

The definition does not itself assert that a realization exists.
In particular, an assertion of the form
\[
\exists\eta\,\Real_T(\eta)
\]
is intended to establish the existence of a class satisfying the
membership condition determined by the fixed parameters of \(T\).
For this purpose, \(\eta\) must not occur free in \(T\), since
otherwise the existential quantifier would also bind occurrences
belonging to those parameters.

The Class-Term Extension Metatheorem will establish that every
admitted class term has a realization, and that any two realizations
of the same class term are equal.

\nbgpara{6}{Contextual Elimination}

The primitive language contains no non-variable terms. A class
term therefore cannot itself be replaced by a term of the primitive
language. Instead, it must be eliminated within the complete
formula in which it occurs.

The elimination procedure uses fresh variables to represent
the classes specified by class terms.

\begin{definition}[Contextual Elimination]
\label{def:contextual-elimination}
Before the procedure is applied, every bounded class term is
expanded into its corresponding general class term.

Let \(\varphi\) be an extended formula containing at least one
class-term occurrence. Select the leftmost outermost such
occurrence, say \(T\), where an occurrence is \emph{outermost}
iff it does not lie within the defining formula of another
class-term occurrence.

The selected occurrence belongs to a unique atomic extended
subformula \(\theta\) of \(\varphi\).

Choose a variable \(\eta\) that does not occur in \(\varphi\),
and let \(\theta'\) be obtained from \(\theta\) by replacing
only the selected occurrence of \(T\) by \(\eta\).

Let \(\varphi'\) be obtained from \(\varphi\) by replacing
the selected occurrence of \(\theta\) with
\[
\exists\eta
\bigl(
\Real_T(\eta)\land\theta'
\bigr).
\]

A single step of \emph{contextual elimination} prescribes
\[
\varphi
\defequiv
\varphi'.
\]
The surrounding logical context is otherwise unchanged.
\end{definition}

\begin{example}
An occurrence of \(T \in A\)
is eliminated as
\[
\exists\eta
\bigl(
\Real_T(\eta)\land\eta\in A
\bigr).
\]
\end{example}

\begin{example}
An occurrence of \(A \in T\)
is eliminated as
\[
\exists\eta
\bigl(
\Real_T(\eta)\land A\in\eta
\bigr).
\]
\end{example}

\begin{example}
An occurrence of \(A=T\)
is eliminated as
\[
\exists\eta
\bigl(
\Real_T(\eta)\land A=\eta
\bigr).
\]
\end{example}

Only the selected occurrence of a class term is replaced.
If the same class term occurs more than once, its occurrences
are eliminated separately, using fresh realization variables
for each occurrence.

The procedure is repeated until no class-term occurrence
remains. Every step removes one class-builder occurrence:
the selected outer builder disappears, while the class terms
appearing in its canonical membership condition are precisely
those already occurring within its defining formula. Hence
the total number of class-builder occurrences decreases by one
at each step, and the procedure terminates after finitely
many steps.

The resulting expression contains no class terms. All remaining
defined notation can then be eliminated according to the
primitive-expansion clauses established previously.

\begin{definition}[Primitive Expansion of a Formula Containing Class Terms]
\label{def:primitive-expansion-with-class-terms}
A \emph{primitive expansion} of an extended formula containing
class terms is obtained by repeatedly applying contextual
elimination until no class terms remain, and then expanding
all remaining defined notation until a formula of the
primitive language is obtained.
\end{definition}

The procedure just specified applies to the extended formulas
generated by the class-term formation clauses, together with
the previously admitted logical and defined notation. The
application of earlier defined predicates and relations to
arbitrary class terms, where not yet licensed by their formation
clauses, will be justified separately through extended-term
substitution.

Although termination follows directly from the construction,
further properties require proof. In particular, we must show
that legitimate choices of fresh variables yield primitive
expansions differing only by alphabetic variation, and that
the procedure respects free variables and capture-avoiding
substitution.

These properties will be established by the Class-Term Extension
Metatheorem.

Contextual elimination must also be distinguished from
predicative reduction. The formula \(\Real_T(\eta)\) contains
an unrestricted universal quantifier, and contextual
elimination introduces an unrestricted existential quantifier.
Consequently, primitive expansion does not by itself establish
that a formula containing class terms may occur within a
predicative defining condition.

Separate predicative reductions will be required for that
purpose. Unlike contextual elimination, these reductions
will use only syntax admitted by Class Formation and will
replace atomic formulas containing class terms by equivalent
formulas involving terms of lower builder rank.

Together, contextual elimination and predicative reduction
will establish the two required aspects of class-builder
notation: its complete eliminability in favor of the
primitive language and its admissibility within the
predicative conditions used to form further classes.

\subsubsection{Class-Builder Metatheorem}

\nbgpara{7}{A Membership Abbreviation}

To formulate the predicative reductions uniformly, we introduce
an auxiliary metalinguistic abbreviation.

Let \(T\) be an extended term and let \(\zeta\) be a variable.
We write
\[
M_T(\zeta)
\]
for the following extended formula:
\[
M_T(\zeta)
\defequiv
\begin{cases}
\zeta\in T,
&
\text{if \(T\) is a variable},\\[2mm]
\operatorname{Cond}_T(\zeta),
&
\text{if \(T\) is a class term}.
\end{cases}
\]

In the second case, the bound variable of \(T\) is renamed
alphabetically whenever necessary to make the canonical membership
condition legitimate.

The abbreviation \(M_T(\zeta)\) expresses the membership condition
associated with \(T\), without presupposing that arbitrary extended
terms may already be substituted for variables.

Two properties will be useful.

First, if \(T\) is a class term of builder rank \(n+1\), then
\(M_T(\zeta)\) contains class terms only of rank at most \(n\).
If \(T\) is a variable, \(M_T(\zeta)\) contains no class terms.

Second, in either case,
\[
\vdash
M_T(\zeta)\to\Set(\zeta).
\]
For variable terms, this follows from Set-I. For general class
terms, it follows immediately from the definition of the
canonical membership condition. The bounded case follows through
its expansion into a general class term.

These observations concern only formulas containing terms already
admitted by the construction. They do not assert a general
substitution principle for extended terms.

\nbgpara{8}{The Class-Term Extension Metatheorem}

The definitions of the preceding paragraphs specify a stagewise
extension of the syntax, but they do not yet establish that the
construction can be continued through every finite builder rank.

In particular, the formation of a class term requires its defining
condition to be predicative. We must therefore show that formulas
containing previously admitted class terms have the predicative
reductions required for their use in further definitions.

At the same time, we must establish that each admitted term has a
unique realization and that contextual elimination yields
well-defined primitive expansions.

These properties will be proved together.

\begin{metatheorem}[Class-Term Extension]
\label{meta:class-term-extension}

The extension by general and bounded class terms is well defined
at every finite builder rank. More precisely:

\begin{enumerate}

\item Every admitted class term \(T\) has a realization.
If \(\eta\) does not occur free in \(T\), then
\[
\vdash
\exists\eta\,\Real_T(\eta).
\]

Any two realizations of \(T\) are equal. Thus, whenever the
displayed formulas are defined,
\[
\vdash
\bigl(
\Real_T(\eta)\land\Real_T(\zeta)
\bigr)
\to
\eta=\zeta.
\]

\item Repeated contextual elimination terminates and assigns every
extended formula containing class terms a primitive expansion,
determined up to alphabetic variation.

The expansion preserves the free variables of the original formula.
It is compatible with capture-avoiding substitution of variables
for variables, provided that freshly introduced bound variables
are renamed where necessary.

\item Every atomic extended formula containing a class term admits
a justified predicative reduction to an extended formula whose
class terms, if any, have strictly lower builder rank than the
greatest builder rank occurring in the original atomic formula.

The reduction introduces no new free variables.

\item Repeated application of these atomic reductions terminates.
Consequently, every formula constructed from the admitted
predicative syntax and class terms can be reduced to a
class-term-free predicative extended formula, and thence to a
core predicative formula by the previously justified reductions.

The resulting reduction is justified independently of any
application of Class Formation whose availability depends upon
that reduction.

\end{enumerate}
\end{metatheorem}

\begin{proof}
We proceed by induction on builder rank.

At rank \(0\), the extended terms are variables. The applicable
formation rules, primitive-expansion procedures, and predicative
reductions have already been established before the introduction
of class-builder notation. Thus, the required properties hold
at the initial stage.

Suppose that the construction has been justified through rank
\(n\). In particular, assume that the class terms of rank at
most \(n\) have unique realizations, that formulas containing
them admit primitive expansions with the stated syntactic
properties, and that the prescribed predicative reductions
have been justified for atomic formulas containing those terms.

We establish the corresponding results at rank \(n+1\).

\medskip

\noindent\textit{Existence and uniqueness of realizations.}

Let
\[
T=\{\xi\mid\varphi\}
\]
be an admitted general class term of rank \(n+1\).

By the formation condition for class terms, \(\varphi\) is
predicative according to the reductions already justified
at lower ranks. Moreover, every class term occurring in
\(\varphi\) has rank at most \(n\).

Choose variables \(\alpha\) and \(\eta\) not occurring in \(T\),
with \(\alpha\neq\eta\).

By the induction hypothesis, capture-avoiding substitution of
\(\alpha\) for \(\xi\) preserves the predicative admissibility
of \(\varphi\). Consequently,
\[
\operatorname{Cond}_T(\alpha)
\defequiv
\Set(\alpha)\land\varphi[\alpha/\xi]
\]
is predicative.

The conditions of Class Formation by Predicative Reduction
are therefore satisfied. Hence
\[
\vdash
\exists\eta\,\forall_S\alpha
\bigl(
\alpha\in\eta
\leftrightarrow
\operatorname{Cond}_T(\alpha)
\bigr).
\]

By the equivalence between restricted and unrestricted
membership specifications, this yields
\[
\vdash
\exists\eta\,\forall\alpha
\left(
\alpha\in\eta
\leftrightarrow
\bigl(
\Set(\alpha)\land
\operatorname{Cond}_T(\alpha)
\bigr)
\right).
\]

Since
\[
\vdash
\operatorname{Cond}_T(\alpha)
\to
\Set(\alpha),
\]
propositional reasoning simplifies the latter specification to
\[
\vdash
\exists\eta\,\forall\alpha
\bigl(
\alpha\in\eta
\leftrightarrow
\operatorname{Cond}_T(\alpha)
\bigr).
\]

By the definition of realization, this is
\[
\vdash
\exists\eta\,\Real_T(\eta).
\]

Thus, \(T\) has a realization.

For uniqueness, suppose that
\[
\Real_T(\eta)
\qquad\text{and}\qquad
\Real_T(\zeta).
\]
By their definitions,
\[
\forall\alpha
\bigl(
\alpha\in\eta
\leftrightarrow
\operatorname{Cond}_T(\alpha)
\bigr)
\]
and
\[
\forall\alpha
\bigl(
\alpha\in\zeta
\leftrightarrow
\operatorname{Cond}_T(\alpha)
\bigr).
\]
It follows that
\[
\forall\alpha
\bigl(
\alpha\in\eta
\leftrightarrow
\alpha\in\zeta
\bigr).
\]
Extensionality therefore gives
\[
\eta=\zeta.
\]

Existence and uniqueness for bounded class terms follow through
their prescribed expansions into general class terms.

\medskip

\noindent\textit{Contextual elimination and primitive expansion.}

Consider a formula whose class terms have rank at most \(n+1\).

A single step of contextual elimination replaces a selected
class-term occurrence \(T\) by a fresh variable \(\eta\) and
introduces the formula
\[
\Real_T(\eta).
\]

The selected outer builder is thereby removed. The class terms
occurring within \(\Real_T(\eta)\) are precisely those inherited
from the defining formula of \(T\), and hence have strictly
lower rank than \(T\).

As already observed in the definition of contextual elimination,
each step decreases the total number of class-builder occurrences
by one. The procedure therefore terminates after finitely many
steps.

When no class terms remain, all remaining defined notation can
be eliminated by the previously established expansion clauses.
The result is a formula of the primitive language.

The choice of a fresh realization variable does not affect this
result except by alphabetic variation. Indeed, two legitimate
choices introduce bound variables in corresponding positions,
and the resulting formulas are immediate alphabetic variants
after any necessary renaming of bound variables already
occurring in the surrounding expression.

Applying this observation successively to the finite elimination
procedure shows that the resulting primitive expansion is
determined up to alphabetic variation.

The procedure also preserves free variables. Replacing \(T\)
by \(\eta\) removes the free variables contributed by that
occurrence of \(T\), but the formula \(\Real_T(\eta)\)
reintroduces precisely those variables as free parameters.
The new variable \(\eta\) is bound by the introduced existential
quantifier. Thus, no free variable is gained or lost.

Compatibility with substitution of variables for variables
follows by induction on the formation of extended expressions.
In the atomic cases, contextual elimination commutes with
substitution after choosing its realization variables fresh
for both the original expression and the variable being
substituted. The logical connectives preserve this property
immediately. At quantifiers and class-builder binders, bound
variables are renamed alphabetically whenever necessary to
avoid capture.

The result is that substitution performed before contextual
elimination and substitution performed after contextual
elimination yield alphabetically variant primitive formulas,
provided that the fresh bound variables in the latter
expansion have first been chosen suitably.

The same argument establishes the corresponding compatibility
of alphabetic variation.

Hence the primitive-expansion procedure has the required
syntactic properties at rank \(n+1\).

\medskip

\noindent\textit{Membership and identity with variable arguments.}

We next establish two equivalences needed to justify the
predicative reductions.

Let \(T\) be a class term whose existence and uniqueness
have just been established, and let \(\zeta\) be a variable
for which its canonical membership condition is defined.

By contextual elimination,
\[
\zeta\in T
\defequiv
\exists\eta
\bigl(
\Real_T(\eta)\land\zeta\in\eta
\bigr),
\]
where \(\eta\) is fresh.

If \(\zeta\in T\), existential elimination supplies a
realization \(\eta\) with \(\zeta\in\eta\). The realization
condition gives
\[
\zeta\in\eta
\leftrightarrow
\operatorname{Cond}_T(\zeta),
\]
and hence \(\operatorname{Cond}_T(\zeta)\).

Conversely, suppose that
\[
\operatorname{Cond}_T(\zeta).
\]
Since a realization of \(T\) exists, existential elimination
supplies a variable \(\eta\) satisfying \(\Real_T(\eta)\).
The same membership equivalence yields
\[
\zeta\in\eta.
\]
Therefore,
\[
\exists\eta
\bigl(
\Real_T(\eta)\land\zeta\in\eta
\bigr),
\]
which is the contextual expansion of \(\zeta\in T\).

We have thus established
\[
\vdash
\zeta\in T
\leftrightarrow
\operatorname{Cond}_T(\zeta).
\]

Similarly, contextual elimination gives
\[
\zeta=T
\defequiv
\exists\eta
\bigl(
\Real_T(\eta)\land\zeta=\eta
\bigr).
\]

From this existential formula, identity elimination yields
\(\Real_T(\zeta)\). Conversely, from \(\Real_T(\zeta)\),
identity reflexivity and existential introduction yield the
displayed existential formula.

Consequently,
\[
\vdash
\zeta=T
\leftrightarrow
\Real_T(\zeta).
\]

Both arguments use only variable-based instances of the
previously established inference rules. They do not presuppose
the substitution of class terms for variables.

The resulting equivalences will be stated separately as the
fundamental Membership and Identity Laws below.

\medskip

\noindent\textit{Predicative reductions of atomic formulas.}

We now assign reductions to atomic extended formulas containing
class terms.

There are three cases. Throughout, variables introduced by a
reduction are chosen distinct from one another and absent from
the surrounding formula.

\smallskip

\noindent\emph{Case 1. Membership with a variable on the left.}

Let \(\zeta\) be a variable and \(T\) a class term.

The membership equivalence just established gives
\[
\vdash
\zeta\in T
\leftrightarrow
M_T(\zeta).
\]

We therefore prescribe
\[
\boxed{
\zeta\in T
\quad\rightsquigarrow\quad
M_T(\zeta).
}
\]

Since \(M_T(\zeta)\) contains only class terms of lower rank
than \(T\), this reduction strictly lowers builder rank.

\smallskip

\noindent\emph{Case 2. Membership with a class term on the left.}

Let \(S\) be a class term and \(T\) any extended term.

Choose fresh variables \(h\) and \(a\).

Contextual elimination of the occurrence of \(S\), followed
by the variable-argument membership equivalence for \(T\),
gives
\[
\vdash
S\in T
\leftrightarrow
\exists h
\bigl(
\Real_S(h)\land M_T(h)
\bigr).
\]

If \(T\) is a variable, the replacement of \(h\in T\) by
\(M_T(h)\) is simply an instance of the definition of \(M\).
If \(T\) is a class term, the replacement is justified by
the membership equivalence established above.

Now
\[
\vdash M_T(h)\to\Set(h).
\]
Furthermore, by the equivalence between restricted and
unrestricted membership specifications and the fact that
\(M_S(a)\) implies \(\Set(a)\), the realization condition
\[
\Real_S(h)
\]
is equivalent to
\[
\forall_S a
\bigl(
a\in h\leftrightarrow M_S(a)
\bigr).
\]

It follows that
\[
\vdash
S\in T
\leftrightarrow
\exists_S h
\left(
\forall_S a
\bigl(
a\in h\leftrightarrow M_S(a)
\bigr)
\land M_T(h)
\right).
\]

We therefore prescribe
\[
\boxed{
S\in T
\quad\rightsquigarrow\quad
\exists_S h
\left(
\forall_S a
\bigl(
a\in h\leftrightarrow M_S(a)
\bigr)
\land M_T(h)
\right).
}
\]

The reduction uses only restricted quantifiers. Every
class term occurring on the right belongs to the defining
formula of \(S\) or \(T\), and therefore has strictly lower
rank than the corresponding outer class term.

\smallskip

\noindent\emph{Case 3. Identity of extended terms.}

Let \(S\) and \(T\) be extended terms, at least one of which
is a class term.

Choose a fresh variable \(a\).

We claim that
\[
\vdash
S=T
\leftrightarrow
\forall_S a
\bigl(
M_S(a)\leftrightarrow M_T(a)
\bigr).
\]

To establish this, consider the realizations of any class
terms occurring in \(S\) and \(T\). Where one of these terms
is a variable, that variable itself is used.

By the variable-argument membership equivalence, the
members of each represented class are precisely the objects
satisfying its associated formula \(M_S(a)\) or \(M_T(a)\).

If the represented classes are equal, identity elimination
therefore gives agreement of their membership conditions.
Universal introduction, followed by restriction to sets,
establishes
\[
\forall_S a
\bigl(
M_S(a)\leftrightarrow M_T(a)
\bigr).
\]

Conversely, suppose that this restricted universal formula
holds.

Every member of any class is a set, and both \(M_S(a)\) and
\(M_T(a)\) imply \(\Set(a)\). Thus, the restricted equivalence
establishes agreement of the two membership conditions for
all objects.

Extensionality gives equality of the corresponding
realizations. Using the existence of realizations and
contextual elimination, we obtain \(S=T\).

Hence the claimed equivalence holds, and we prescribe
\[
\boxed{
S=T
\quad\rightsquigarrow\quad
\forall_S a
\bigl(
M_S(a)\leftrightarrow M_T(a)
\bigr).
}
\]

Again, every class term occurring on the right has strictly
lower rank than an outer class term occurring on the left.

\medskip

These three cases exhaust the atomic extended formulas
containing class terms. The only atomic relation symbols are
\(\in\) and \(=\), and the displayed cases cover every
possible placement of class terms in their arguments.

All three reductions are justified by derivable equivalences.
They introduce no free variables, since their newly
introduced variables are bound and chosen fresh for the
surrounding formula.

They also introduce no unrestricted quantifiers. Their
right-hand sides are constructed using the atomic and
previously admitted predicative notation, logical
connectives, restricted quantifiers, and class terms of
lower rank.

\medskip

\noindent\textit{Termination and predicative admissibility.}

It remains to establish that the prescribed reductions
can be applied repeatedly without circularity.

Consider the multiset of builder ranks of all class-term
occurrences in a finite extended formula.

Every atomic reduction removes the outer class terms
occurring in the selected atomic formula and replaces them
with formulas containing only their lower-ranked constituent
terms. Consequently, each application strictly decreases
the multiset of builder ranks under the standard multiset
ordering.

Since this ordering is well founded on finite multisets
of natural numbers, repeated reduction terminates.

More concretely, one may reduce all atomic occurrences
containing terms of the greatest rank, then those
containing terms of the next lower rank, and so on.
No reduction introduces a term of the rank just eliminated
or of any higher rank.

When no class terms remain, the resulting formula belongs
to the previously established extended syntax.

If the original formula was constructed using only the
admitted predicative operations, its class-term-free
reduct is predicative according to the reductions already
justified before class-builder notation was introduced.
Repeated application of those reductions therefore yields
a core predicative formula.

The equivalence of the original formula and its reduct
follows by successive applications of the atomic
equivalences, together with the established congruence
principles for the logical connectives and quantifiers.

All variable choices are made fresh for the surrounding
formula, so no reduction introduces additional free
variables or captures existing free occurrences. The
prescribed order of reduction determines the result up
to alphabetic variation.

Finally, the construction is non-circular. At rank
\(n+1\), the existence of realizations was established
using only predicative reductions justified at ranks at
most \(n\). The membership and identity equivalences
were then derived from those realizations. Only afterward
were the new atomic reductions justified.

Thus, no application of Class Formation depends upon
a reduction whose justification presupposes that same
application.

The induction step is complete. By induction, the stated
properties hold at every finite builder rank.
\end{proof}